% ============================================================
%  MAG7 TRACKER — Documentación Técnica
%  Para actualizar: editar este .tex y correr build.bat
% ============================================================
\documentclass[11pt, a4paper]{article}

% ── Packages ──────────────────────────────────────────────
\usepackage[utf8]{inputenc}
\usepackage[T1]{fontenc}
\usepackage[spanish]{babel}
\usepackage{lmodern}
\usepackage[margin=2.5cm]{geometry}
\usepackage{amsmath, amssymb}
\DeclareMathOperator*{\argmax}{arg\,max}
\usepackage{booktabs}
\usepackage{array}
\usepackage{longtable}
\usepackage{xcolor}
\usepackage{hyperref}
\usepackage{enumitem}
\usepackage{titlesec}
\usepackage{fancyhdr}
\usepackage{graphicx}
\usepackage{tcolorbox}
\tcbuselibrary{skins, breakable}
\usepackage{multirow}
\usepackage{tabularx}

% ── Colores ───────────────────────────────────────────────
\definecolor{gold}{HTML}{B07800}
\definecolor{pink}{HTML}{C0304A}
\definecolor{dark}{HTML}{1A1A2E}
\definecolor{green}{HTML}{1A7A3C}
\definecolor{red}{HTML}{C0392B}
\definecolor{blue}{HTML}{1A5FA8}
\definecolor{gray}{HTML}{555566}
\definecolor{lightgray}{HTML}{CCCCDD}

% ── Estilo de secciones ───────────────────────────────────
\titleformat{\section}{\large\bfseries\color{pink}}{}{0em}{}[\titlerule]
\titleformat{\subsection}{\normalsize\bfseries\color{gold}}{}{0em}{}
\titleformat{\subsubsection}{\normalsize\bfseries}{}{0em}{}

% ── Header / Footer ───────────────────────────────────────
\pagestyle{fancy}
\fancyhf{}
\fancyhead[L]{\small\textcolor{gray}{MAG7 TRACKER · Documentación Técnica}}
\fancyhead[R]{\small\textcolor{gray}{\thepage}}
\renewcommand{\headrulewidth}{0.4pt}
\renewcommand{\headrule}{\hbox to\headwidth{\color{lightgray}\leaders\hrule height \headrulewidth\hfill}}

% ── Cajas de fórmulas ─────────────────────────────────────
\newtcolorbox{formulabox}[1][]{
  enhanced, breakable,
  colback=gold!10, colframe=gold!70,
  coltitle=brown!60!black,
  fonttitle=\small\bfseries,
  title={#1}, arc=4pt, boxrule=0.6pt,
  left=8pt, right=8pt, top=6pt, bottom=6pt
}

\newtcolorbox{infobox}[1][]{
  enhanced, breakable,
  colback=blue!7, colframe=blue!50,
  coltitle=blue!60!black,
  fonttitle=\small\bfseries,
  title={#1}, arc=4pt, boxrule=0.6pt,
  left=8pt, right=8pt, top=6pt, bottom=6pt
}

\newtcolorbox{warningbox}{
  enhanced, colback=red!8, colframe=red!50,
  arc=4pt, boxrule=0.6pt,
  left=8pt, right=8pt, top=6pt, bottom=6pt
}

% ── Hiperenlaces ──────────────────────────────────────────
\hypersetup{
  colorlinks=true, linkcolor=gold, urlcolor=pink, citecolor=green,
  pdftitle={MAG7 Tracker - Documentacion Tecnica},
  pdfauthor={LarrainVial}
}

% ── Helpers ───────────────────────────────────────────────
\newcommand{\tab}[1]{\textcolor{gold}{\textbf{[#1]}}}
\newcommand{\api}[1]{\texttt{\textcolor{pink}{#1}}}
\newcommand{\code}[1]{\texttt{#1}}

% ============================================================
\begin{document}

% ── Portada ───────────────────────────────────────────────
\begin{titlepage}
  \centering
  \vspace*{3cm}
  {\fontsize{36}{44}\selectfont\bfseries\textcolor{pink}{MAG 7}\par}
  {\fontsize{36}{44}\selectfont\bfseries\textcolor{gold}{TRACKER}\par}
  \vspace{0.8cm}
  {\large\textcolor{gray}{Documentación Técnica}\par}
  \vspace{0.4cm}
  {\normalsize\textcolor{gray}{Magnificent Seven · Performance Analysis Platform}\par}
  \vspace{2cm}
  \rule{8cm}{0.4pt}
  \vspace{1cm}

  \begin{tabular}{ll}
    \textcolor{gray}{Stack}    & FastAPI · React · TypeScript · Recharts \\[4pt]
    \textcolor{gray}{Datos}    & Yahoo Finance API (tiempo real) \\[4pt]
    \textcolor{gray}{Tickers}  & AAPL · MSFT · GOOGL · AMZN · NVDA · META · TSLA \\[4pt]
    \textcolor{gray}{Versión}  & 2.0 \\
  \end{tabular}

  \vfill
  {\small\textcolor{gray}{Para actualizar este documento, editar \texttt{mag7\_doc.tex} y ejecutar \texttt{build.bat}}}
\end{titlepage}

% ── Tabla de contenidos ───────────────────────────────────
\tableofcontents
\newpage

% ============================================================
\section{Arquitectura General}

El MAG7 Tracker es una aplicación web de análisis financiero para las siete empresas tecnológicas más grandes del mundo (\textit{Magnificent Seven}). Está compuesta por:

\begin{itemize}[leftmargin=1.5cm]
  \item \textbf{Backend}: FastAPI (Python) corriendo en \code{localhost:8000}. Consume datos de Yahoo Finance, calcula todos los indicadores y métricas, y expone una REST API.
  \item \textbf{Frontend}: React + TypeScript con Recharts para visualización. Corre en \code{localhost:5173} (Vite).
  \item \textbf{Proxy}: Vite redirige \code{/api/*} al backend, evitando problemas de CORS.
\end{itemize}

\subsection{Endpoints de la API}

\begin{tabularx}{\textwidth}{lX}
  \toprule
  \textbf{Endpoint} & \textbf{Descripción} \\
  \midrule
  \api{GET /api/companies}              & Lista de las 7 empresas con nombre y color \\
  \api{GET /api/prices}                 & Precios históricos normalizados (Base 100 y Vol-Adj) \\
  \api{GET /api/forecast}               & Pronóstico 30 días con 3 modelos estocásticos \\
  \api{GET /api/technical}              & Indicadores técnicos + señales automáticas \\
  \api{GET /api/analytics}              & Métricas de riesgo/retorno comparativas \\
  \api{GET /api/fundamentals}           & Ratios de valoración (P/E, P/B, ROE, etc.) \\
  \api{GET /api/backtest}               & Backtest de estrategias técnicas \\
  \api{GET /api/portfolio}              & Análisis de portafolio + Tangent Portfolio \\
  \api{GET /api/sectors}                & Comparativa por sectores del S\&P 500 \\
  \api{GET /api/news}                   & Noticias recientes de Yahoo Finance \\
  \api{GET /api/ai\_summary}            & Resumen cuantitativo automático (técnico + fundamental) \\
  \bottomrule
\end{tabularx}

\newpage

% ============================================================
\section{Tab 1: Comparison}

Compara la evolución histórica de los Mag7 en un mismo gráfico. Permite seleccionar empresas, rango de fechas y modo de visualización.

\subsection{Base 100}

Todos los precios se normalizan al valor 100 en la fecha de inicio:

\begin{formulabox}[Normalización Base 100]
  \[ P_t^{(100)} = 100 \cdot \frac{P_t}{P_{t_0}} \]
  Donde $P_{t_0}$ es el precio en el primer día del período. Permite comparar retornos independientemente del precio absoluto.
\end{formulabox}

\subsection{Vol-Adjusted}

Divide el retorno acumulado por la volatilidad anualizada, similar a normalizar por el riesgo asumido:

\begin{formulabox}[Volatilidad Ajustada]
  \[ P_t^{(\text{vol-adj})} = \frac{P_t^{(100)}}{\hat{\sigma} \cdot \sqrt{252}} \]
  Donde $\hat{\sigma}$ es la desviación estándar de los log-retornos diarios. Penaliza a acciones muy volátiles.
\end{formulabox}

\subsection{Modelos de Pronóstico a 30 días}

Al activar el forecast se proyectan 30 días hacia adelante con bandas de confianza. Se pueden elegir tres modelos:

\subsubsection{GBM — Geometric Brownian Motion}

\begin{formulabox}[GBM]
  \[ S_{t+1} = S_t \cdot \exp\!\left[\left(\mu - \frac{\sigma^2}{2}\right)\Delta t + \sigma \sqrt{\Delta t}\, Z\right], \quad Z \sim \mathcal{N}(0,1) \]
  \begin{itemize}
    \item $\mu$: drift (media de log-retornos anualizados)
    \item $\sigma$: volatilidad histórica anualizada
    \item \textbf{Pro}: Simple, analítico, base del modelo Black-Scholes
    \item \textbf{Con}: Asume volatilidad constante, ignora \textit{fat tails}
  \end{itemize}
\end{formulabox}

\subsubsection{GARCH(1,1) — Volatilidad Condicional}

\begin{formulabox}[{GARCH(1,1)}]
  \[ \sigma_t^2 = \omega + \alpha \cdot \varepsilon_{t-1}^2 + \beta \cdot \sigma_{t-1}^2 \]
  \begin{itemize}
    \item La varianza de mañana depende del shock de hoy $(\varepsilon^2)$ y la varianza de ayer $(\sigma^2)$
    \item Captura \textit{volatility clustering}: períodos de alta vol seguidos de alta vol
    \item Parámetros $\{\omega, \alpha, \beta\}$ estimados por máxima verosimilitud
    \item \textbf{Pro}: Más realista para series financieras
    \item \textbf{Con}: No modela saltos bruscos
  \end{itemize}
\end{formulabox}

\subsubsection{Merton Jump-Diffusion}

\begin{formulabox}[Merton Jump-Diffusion]
  \[ \frac{dS}{S} = (\mu - \lambda\kappa)\,dt + \sigma\,dW + J\,dN(\lambda) \]
  \begin{itemize}
    \item $dN(\lambda)$: proceso de Poisson con intensidad $\lambda$ (frecuencia de saltos)
    \item $J \sim \mathcal{N}(\mu_J, \sigma_J^2)$: tamaño del salto en log-escala
    \item Retornos con $|r| > 2.5\sigma$ se clasifican como ``saltos'' y se modelan por separado
    \item \textbf{Pro}: Captura eventos extremos (\textit{fat tails}, crashes)
    \item \textbf{Con}: Sensible al umbral de clasificación de saltos
  \end{itemize}
\end{formulabox}

Las bandas de confianza se construyen con 1000 simulaciones Monte Carlo; se reportan los percentiles 10 y 90 (para CI=80\%).

\newpage

% ============================================================
\section{Tab 2: Technical}

Análisis técnico individual de cada stock. Incluye indicadores clásicos, señales automáticas, backtesting y noticias.

\subsection{Indicadores de Precio}

\subsubsection{SMA — Simple Moving Average}

\begin{formulabox}[{SMA(n)}]
  \[ \text{SMA}_t(n) = \frac{1}{n} \sum_{i=0}^{n-1} P_{t-i} \]
  Se calculan tres períodos: \textbf{20 días} (corto), \textbf{50 días} (medio), \textbf{200 días} (largo plazo). Suavizan el ruido y revelan la tendencia subyacente.
\end{formulabox}

\subsubsection{Bollinger Bands — BB(20, $\pm 2\sigma$)}

\begin{formulabox}[Bollinger Bands]
  \begin{align}
    \text{BB}^{\text{mid}}_t  &= \text{SMA}_t(20) \\
    \text{BB}^{\text{upper}}_t &= \text{BB}^{\text{mid}}_t + 2 \cdot \hat{\sigma}_t(20) \\
    \text{BB}^{\text{lower}}_t &= \text{BB}^{\text{mid}}_t - 2 \cdot \hat{\sigma}_t(20)
  \end{align}
  El 95\% de los precios debería quedar dentro de las bandas bajo distribución normal. Precio en banda superior = posible sobrecompra; en banda inferior = posible sobreventa.
\end{formulabox}

\subsection{Indicadores de Momentum}

\subsubsection{RSI(14) — Relative Strength Index}

\begin{formulabox}[RSI de Wilder]
  \[ \text{RSI}_t = 100 - \frac{100}{1 + RS_t}, \quad RS_t = \frac{\overline{G}_{14}}{\overline{L}_{14}} \]
  Donde $\overline{G}_{14}$ y $\overline{L}_{14}$ son la media exponencial de ganancias y pérdidas de 14 días (EMA de Wilder: $\alpha = 1/14$).
  \begin{itemize}
    \item RSI $> 70$: \textbf{sobrecompra} — momentum agotado
    \item RSI $< 30$: \textbf{sobreventa} — posible rebote
    \item RSI = 50: zona neutral
  \end{itemize}
\end{formulabox}

\subsubsection{MACD(12, 26, 9) — Moving Average Convergence Divergence}

\begin{formulabox}[MACD]
  \begin{align}
    \text{MACD}_t      &= \text{EMA}_{12}(P) - \text{EMA}_{26}(P) \\
    \text{Signal}_t    &= \text{EMA}_9(\text{MACD}) \\
    \text{Hist}_t      &= \text{MACD}_t - \text{Signal}_t
  \end{align}
  \begin{itemize}
    \item \textbf{Cruce alcista}: MACD cruza sobre la señal → impulso positivo
    \item \textbf{Cruce bajista}: MACD cruza bajo la señal → impulso negativo
    \item \textbf{Histograma verde}: momentum alcista creciente
    \item \textbf{Histograma rojo}: momentum bajista creciente
  \end{itemize}
\end{formulabox}

\subsection{Earnings Markers}

Los resultados trimestrales (EPS) se marcan en los gráficos con un triángulo dorado (\textbf{E}). Al hacer hover se despliega:

\begin{center}
\begin{tabular}{ll}
  \toprule
  \textbf{Campo} & \textbf{Descripción} \\
  \midrule
  Date           & Fecha del reporte \\
  EPS Estimate   & Consenso de analistas (dólares por acción) \\
  EPS Actual     & EPS reportado por la empresa \\
  EPS Surprise   & $\frac{\text{Actual} - \text{Estimate}}{\text{Estimate}} \times 100\%$ \\
  \bottomrule
\end{tabular}
\end{center}

Verde = superó estimados; Rojo = no alcanzó estimados.

\subsection{Señales Técnicas Automáticas}

Al cargar el análisis, el sistema detecta automáticamente señales basadas en los indicadores más recientes:

\begin{center}
\begin{tabularx}{\textwidth}{llX}
  \toprule
  \textbf{Indicador} & \textbf{Condición} & \textbf{Señal} \\
  \midrule
  RSI & $< 30$ & \textcolor{green}{$\blacktriangle$} Bullish fuerte: sobreventa \\
  RSI & $30$--$40$ & \textcolor{gold}{$\blacktriangleright$} Bullish débil: momentum bajo \\
  RSI & $> 70$ & \textcolor{red}{$\blacktriangledown$} Bearish fuerte: sobrecompra \\
  RSI & $60$--$70$ & \textcolor{gold}{$\blacktriangleright$} Bearish débil: cerca de sobrecompra \\
  \midrule
  MACD & Cruza sobre señal & \textcolor{green}{$\blacktriangle$} Bullish fuerte: cruce alcista \\
  MACD & Cruza bajo señal  & \textcolor{red}{$\blacktriangledown$} Bearish fuerte: cruce bajista \\
  MACD & Por encima señal  & \textcolor{gold}{$\blacktriangleright$} Bullish débil: momentum positivo \\
  MACD & Por debajo señal  & \textcolor{gold}{$\blacktriangleright$} Bearish débil: momentum negativo \\
  \midrule
  SMA 50/200 & SMA50 cruza $\uparrow$ SMA200 & \textcolor{green}{$\bigstar$} Bullish fuerte: \textit{Golden Cross} \\
  SMA 50/200 & SMA50 cruza $\downarrow$ SMA200 & \textcolor{red}{$\times$} Bearish fuerte: \textit{Death Cross} \\
  \midrule
  BB & Precio toca banda superior & \textcolor{red}{$\blacktriangledown$} Bearish fuerte: posible reversión \\
  BB & Precio toca banda inferior & \textcolor{green}{$\blacktriangle$} Bullish fuerte: posible rebote \\
  \midrule
  Tendencia & Precio $> 10\%$ bajo SMA200 & \textcolor{red}{$\Downarrow$} Bearish fuerte: tendencia bajista \\
  Tendencia & Precio $> 25\%$ sobre SMA200 & \textcolor{gold}{$\Uparrow$} Bearish débil: muy extendido \\
  \bottomrule
\end{tabularx}
\end{center}

El resumen (Alcista/Bajista/Neutral) se calcula contando señales fuertes (peso 2) y débiles (peso 1).

\subsection{Backtesting de Estrategias}

Permite testear cuatro estrategias técnicas sobre el período histórico seleccionado:

\begin{center}
\begin{tabularx}{\textwidth}{llX}
  \toprule
  \textbf{Estrategia} & \textbf{Entrada} & \textbf{Salida} \\
  \midrule
  RSI Mean Revert  & RSI cruza bajo 30  & RSI cruza sobre 70 \\
  MACD Crossover   & MACD $>$ señal     & MACD $<$ señal \\
  SMA Golden/Death & SMA50 $>$ SMA200   & SMA50 $<$ SMA200 \\
  BB Mean Revert   & Precio $<$ BB lower & Precio $>$ BB middle \\
  \bottomrule
\end{tabularx}
\end{center}

\textbf{Supuestos del backtest:} solo posiciones largas (long/cash), sin apalancamiento, sin comisiones, ejecución al precio de cierre del día de la señal. Curva de equity base 100, comparada contra Buy \& Hold.

\textbf{Métricas reportadas:}
\begin{itemize}
  \item \textbf{Retorno Estrategia vs B\&H}: retorno total de cada enfoque
  \item \textbf{Exceso}: diferencia de retorno entre estrategia y Buy \& Hold
  \item \textbf{Sharpe}: $(\bar{r} \cdot 252 - r_f) / (\hat{\sigma} \cdot \sqrt{252})$, con $r_f = 4.5\%$
  \item \textbf{Max Drawdown}: caída máxima desde el pico de la curva de equity
  \item \textbf{Win Rate}: porcentaje de trades ganadores
  \item \textbf{Tiempo en Mercado}: porcentaje de días con posición abierta
\end{itemize}

\subsection{Resumen Cuantitativo Automático}

Al pulsar \textbf{Resumen IA} el sistema genera un análisis de 4 oraciones en español sin llamadas externas. El endpoint \api{/api/ai\_summary} calcula todos los indicadores en el backend y construye el texto mediante lógica condicional determinista:

\begin{center}
\begin{tabularx}{\textwidth}{llX}
  \toprule
  \textbf{Oración} & \textbf{Fuente} & \textbf{Contenido} \\
  \midrule
  1 — Tendencia  & RSI, MACD, SMA50/200, BB\%  & Diagnóstico técnico general \\
  2 — Fundamentales & P/E, Rev.Growth, Margen, ROE & Calidad del negocio \\
  3 — Riesgos    & RSI $>70$, P/E $>40$, deuda, sorpresa EPS & Alertas y riesgos actuales \\
  4 — Sesgo      & Score ponderado de señales & Conclusión alcista/bajista/neutral \\
  \bottomrule
\end{tabularx}
\end{center}

El resultado se cachea 5 minutos por ticker (\code{\_AI\_SUMMARY\_CACHE}) para evitar recalcular en cada petición.

\newpage

% ============================================================
\section{Tab 3: Analytics}

Análisis cuantitativo comparativo de todos los stocks seleccionados versus el S\&P 500 (SPY) como benchmark. Las secciones siguen una narrativa de valuación progresiva: consenso externo → modelo propio → score cuantitativo → riesgo/retorno → datos raw.

\subsection{Precio Objetivo de Analistas}

Muestra el consenso de precio objetivo de Wall Street para cada stock. Los datos se obtienen del módulo \code{financialData} de Yahoo Finance (\code{/api/fundamentals}, requiere pulsar \textbf{Load Fundamentals}).

\begin{center}
\begin{tabular}{ll}
  \toprule
  \textbf{Campo} & \textbf{Descripción} \\
  \midrule
  Precio Actual       & Último precio de mercado \\
  Objetivo Medio      & Media de los precios objetivo de analistas \\
  Objetivo Máximo     & Estimación más optimista \\
  Objetivo Mínimo     & Estimación más conservadora \\
  Upside / Downside   & $(\text{Medio} - \text{Actual}) / \text{Actual} \times 100\%$ \\
  Recomendación       & Consenso textual (Strong Buy / Buy / Hold / Sell) \\
  N° Analistas        & Número de analistas que cubren el valor \\
  \bottomrule
\end{tabular}
\end{center}

La barra de rango visualiza el espectro mínimo--media--máximo en escala relativa, permitiendo ver la dispersión de opiniones.

\subsection{DCF Interactivo}

Modelo de Descuento de Flujos de Caja (DCF) de dos etapas con supuestos editables en tiempo real. Requiere \textbf{Load Fundamentals}.

\begin{formulabox}[DCF de 2 Etapas — Fórmula General]
\textbf{Proyección de FCF:}
\[
  \text{FCF}_t =
  \begin{cases}
    \text{FCF}_0 \cdot (1 + g_1)^t & t = 1, \ldots, 5 \\
    \text{FCF}_5 \cdot (1 + g_2)^{t-5} & t = 6, \ldots, 10
  \end{cases}
\]
\textbf{Valor Terminal (Gordon Growth):}
\[
  \text{TV} = \frac{\text{FCF}_{10} \cdot (1 + g_T)}{\text{WACC} - g_T}
\]
\textbf{Enterprise Value:}
\[
  \text{EV} = \sum_{t=1}^{10} \frac{\text{FCF}_t}{(1+\text{WACC})^t} + \frac{\text{TV}}{(1+\text{WACC})^{10}}
\]
\textbf{Precio intrínseco por acción:}
\[
  P^* = \frac{\text{EV} - \text{Deuda Neta}}{\text{Acciones en circulación}}
\]
\end{formulabox}

\textbf{Supuestos por defecto} (derivados automáticamente de los datos TTM de Yahoo Finance):

\begin{center}
\begin{tabular}{lll}
  \toprule
  \textbf{Parámetro} & \textbf{Default} & \textbf{Fuente} \\
  \midrule
  FCF Base          & Free Cash Flow TTM           & \code{financialData.freeCashflow} \\
  $g_1$ (años 1--5) & Crecimiento de revenue TTM   & \code{financialData.revenueGrowth} \\
  $g_2$ (años 6--10) & $g_1 / 2$                  & Fade automático \\
  $g_T$ Terminal    & 2.5\%                        & Constante conservadora \\
  WACC              & $r_f + \beta \times 5.5\%$   & CAPM ($r_f = 4.5\%$, ERP = 5.5\%) \\
  Deuda Neta        & Deuda total $-$ Caja         & \code{totalDebt} $-$ \code{totalCash} \\
  Acciones          & Shares outstanding           & \code{defaultKeyStatistics} \\
  \bottomrule
\end{tabular}
\end{center}

El usuario puede ajustar todos los parámetros con sliders. El botón \textbf{Reset} restaura los valores derivados de los datos. El gráfico de barras muestra el valor presente de cada año (amarillo = años 1--5, rosa = años 6--10, violeta = valor terminal).

\begin{warningbox}
  Este modelo es solo orientativo. Los resultados dependen fuertemente de los supuestos de crecimiento y WACC. No constituye asesoramiento de inversión.
\end{warningbox}

\subsection{Score de Inversión}

Cada acción recibe un puntaje de 0 a 100 basado en seis métricas cuantitativas (más un bono opcional de fundamentales):

\begin{center}
\begin{tabular}{llr}
  \toprule
  \textbf{Métrica} & \textbf{Lógica} & \textbf{Máx pts} \\
  \midrule
  Sharpe  & 0 pts si $\leq 0$; 20 pts si $\geq 2$; lineal intermedio & 20 \\
  Alpha   & 0 pts si $\leq -10\%$; 20 pts si $\geq 15\%$ & 20 \\
  Calmar  & 0 pts si $\leq 0$; 15 pts si $\geq 2$ & 15 \\
  Sortino & 0 pts si $\leq 0$; 15 pts si $\geq 2$ & 15 \\
  Max DD  & 15 pts si $\geq -5\%$; 0 pts si $\leq -60\%$ & 15 \\
  Beta    & 15 pts si $0.7 \leq \beta \leq 1.2$ (sweet spot) & 15 \\
  \midrule
  \textit{Fundamentales (opcional)} & P/E + crecimiento + márgenes & \textit{15} \\
  \bottomrule
\end{tabular}
\end{center}

Rating: \textbf{\textcolor{green}{$\blacktriangle$ Buy}} $\geq 60/100$ · \textbf{\textcolor{gold}{$\blacktriangleright$ Hold}} $\geq 38/100$ · \textbf{\textcolor{red}{$\blacktriangledown$ Sell}} $< 38/100$.

\subsection{Métricas de Riesgo/Retorno}

\subsubsection{Retorno Total y Volatilidad Anualizada}

\begin{formulabox}[Retorno y Volatilidad]
  \[ r_{\text{total}} = \frac{P_T - P_0}{P_0} \times 100\%, \qquad \hat{\sigma}_{\text{ann}} = \hat{\sigma}_{\text{diaria}} \cdot \sqrt{252} \]
  Donde $\hat{\sigma}_{\text{diaria}}$ es la desviación estándar de los log-retornos diarios $r_t = \ln(P_t/P_{t-1})$.
\end{formulabox}

\subsubsection{Sharpe Ratio}

\begin{formulabox}[{Sharpe Ratio (con $r_f = 4.5\%$)}]
  \[ \text{Sharpe} = \frac{\mu_{\text{ann}} - r_f}{\sigma_{\text{ann}}} \]
  Mide el retorno extra por unidad de riesgo total. $> 1.5$ = excelente; $0.5$--$1.5$ = aceptable; $< 0.5$ = pobre.
\end{formulabox}

\subsubsection{Beta (CAPM)}

\begin{formulabox}[Beta vs S\&P 500]
  \[ \beta_i = \frac{\text{Cov}(r_i, r_{\text{SPY}})}{\text{Var}(r_{\text{SPY}})} \]
  $\beta > 1$: más volátil que el mercado (amplifica movimientos). $\beta < 1$: más defensivo. $\beta = 1$: se mueve en línea con el mercado.
\end{formulabox}

\subsubsection{Alpha de Jensen}

\begin{formulabox}[{Alpha (CAPM)}]
  \[ \alpha_i = r_i - \left[r_f + \beta_i \cdot (r_m - r_f)\right] \]
  Retorno que \textbf{no explica el mercado}. Alpha positivo = la acción generó retorno extra por encima de lo que predice su nivel de riesgo sistemático. $r_f = 4.5\%$, $r_m = $ retorno de SPY.
\end{formulabox}

\subsubsection{Max Drawdown}

\begin{formulabox}[Maximum Drawdown]
  \[ \text{MDD} = \min_{t \in [0,T]} \left(\frac{P_t}{\max_{\tau \leq t} P_\tau} - 1\right) \times 100\% \]
  Caída máxima desde cualquier pico hasta el valle siguiente. Mide el peor escenario histórico para un inversor que compró en el máximo.
\end{formulabox}

\subsubsection{Calmar Ratio}

\begin{formulabox}[Calmar Ratio]
  \[ \text{Calmar} = \frac{r_{\text{ann}}}{|\text{MDD}|} \]
  Retorno anualizado por unidad de máximo drawdown. $> 1$ = bueno; $> 2$ = excelente.
\end{formulabox}

\subsubsection{Sortino Ratio}

\begin{formulabox}[Sortino Ratio]
  \[ \text{Sortino} = \frac{\mu_{\text{ann}} - r_f}{\sigma_{\text{down}} \cdot \sqrt{252}} \]
  Donde $\sigma_{\text{down}}$ es la desviación estándar de los retornos \textbf{negativos solamente}. Penaliza únicamente la volatilidad bajista, a diferencia del Sharpe que penaliza toda volatilidad.
\end{formulabox}

\subsection{Análisis Automático de Insights}

El sistema genera automáticamente 7 insights en texto:
\begin{enumerate}
  \item Mejor y peor desempeño del período
  \item Mejor ratio riesgo/retorno (Sharpe)
  \item Mayor y menor Alpha vs CAPM
  \item Betas extremos (más agresivo vs más defensivo)
  \item Par más y menos correlacionado (diversificación)
  \item Mayor caída desde máximo (drawdown)
  \item Mejor Sortino (volatilidad bajista)
\end{enumerate}

\subsection{Drawdown y Volatilidad Rodante}

\begin{itemize}
  \item \textbf{Drawdown acumulado}: gráfico temporal de la caída desde máximo para cada stock y SPY
  \item \textbf{Volatilidad rodante 21d}: $\hat{\sigma}_t^{(21)} = \text{std}(r_{t-20}, \ldots, r_t) \cdot \sqrt{252}$, permite ver si la volatilidad aumentó o disminuyó en el tiempo
\end{itemize}

\subsection{Matriz de Correlación}

Correlación de Pearson de los log-retornos diarios entre todos los stocks seleccionados:
\[ \rho_{ij} = \frac{\text{Cov}(r_i, r_j)}{\hat{\sigma}_i \hat{\sigma}_j} \in [-1, 1] \]
Visualizada como heatmap con escala de color: 0 = gris, 1 = dorado. Correlaciones altas implican menor diversificación entre esos dos activos.

\subsection{Fundamentales}

Datos obtenidos de Yahoo Finance (\code{quoteSummary}). Se muestran los principales ratios de valoración y calidad:

\begin{center}
\begin{tabularx}{\textwidth}{llX}
  \toprule
  \textbf{Métrica} & \textbf{Sigla} & \textbf{Interpretación} \\
  \midrule
  P/E Trailing  & trailing\_pe & Precio / Ganancias últ. 12 meses. Bajo = barato \\
  P/E Forward   & forward\_pe  & Precio / Ganancias proyectadas. Expectativa del mercado \\
  P/B           & price\_to\_book & Precio / Valor libro. $> 1$ = mercado valora más que contable \\
  EV/EBITDA     & ev\_to\_ebitda  & Múltiplo de empresa. Comparable entre empresas con distinto leverage \\
  PEG           & peg\_ratio     & P/E / Crecimiento. $< 1$ sugiere subvaloración relativa \\
  Rev. Growth   & revenue\_growth & Crecimiento YoY de ingresos \\
  EPS Growth    & earnings\_growth & Crecimiento YoY de ganancias por acción \\
  Profit Margin & profit\_margin  & Ganancia neta / ingresos \\
  ROE           & roe            & Return on Equity: ganancia neta / patrimonio \\
  Deuda/Capital & debt\_to\_equity & Apalancamiento. Alto = más riesgo financiero \\
  Div. Yield    & dividend\_yield & Dividendo / precio. Relevante para inversión por renta \\
  \bottomrule
\end{tabularx}
\end{center}

\subsection{Radar de Valuación}

Visualización tipo \textit{spider chart} que compara todos los stocks seleccionados en 6 ejes simultáneamente. Cada eje se normaliza al rango $[0, 100]$ usando \textbf{min-max scaling} sobre los tickers disponibles:

\[
  x_{\text{norm}} = \frac{x - \min}{\max - \min} \times 100
\]

\begin{center}
\begin{tabular}{lll}
  \toprule
  \textbf{Eje} & \textbf{Fuente} & \textbf{Dirección favorable} \\
  \midrule
  Crecimiento    & Revenue growth YoY      & Mayor = mejor \\
  Rentabilidad   & Profit margin           & Mayor = mejor \\
  Valuación      & $100 - \text{EV/EBITDA normalizado}$ & Menor múltiplo = mejor \\
  ROE            & Return on Equity        & Mayor = mejor \\
  Momentum EPS   & EPS growth YoY          & Mayor = mejor \\
  Calidad        & $100 - \text{Deuda/Capital normalizado}$ & Menor deuda = mejor \\
  \bottomrule
\end{tabular}
\end{center}

El radar permite identificar de un vistazo qué empresa lidera en cada dimensión y si alguna domina en múltiples frentes. Cada stock se representa con el color corporativo de la empresa.

\newpage

% ============================================================
\section{Tab 4: Portfolio}

Análisis de un portafolio personalizado con los stocks seleccionados, con pesos definidos por el usuario.

\subsection{Portafolio Manual}

El usuario asigna pesos porcentuales $w_i$ a cada acción ($\sum w_i = 100\%$). El retorno del portafolio es:
\[ r_P = \sum_{i=1}^{N} w_i \cdot r_i \]

Se muestran las métricas del portafolio (mismo set que Analytics) comparadas contra SPY.

\subsection{Tangent Portfolio (Markowitz)}

\begin{formulabox}[Portafolio Tangente]
  El portafolio tangente es la combinación de pesos $\mathbf{w}^*$ que \textbf{maximiza el Sharpe ratio}:
  \[ \mathbf{w}^* = \argmax_{\mathbf{w}} \frac{\mathbf{w}^T \boldsymbol{\mu} - r_f}{\sqrt{\mathbf{w}^T \boldsymbol{\Sigma} \mathbf{w}}} \]
  Sujeto a $\mathbf{w} \geq 0$ y $\sum w_i = 1$.\\[6pt]
  Donde $\boldsymbol{\mu}$ es el vector de retornos esperados (históricos anualizados) y $\boldsymbol{\Sigma}$ es la matriz de covarianzas de los retornos diarios, anualizada por $\times 252$.
\end{formulabox}

La optimización se resuelve numéricamente con \code{scipy.optimize.minimize} (método SLSQP). El resultado muestra los pesos óptimos y la curva de equity del portafolio teóricamente óptimo sobre el período histórico.

\begin{infobox}[Nota importante]
  Los pesos del Tangent Portfolio se calculan sobre datos históricos y no garantizan resultados futuros. Es una referencia para entender qué combinación hubiera maximizado el Sharpe en el pasado.
\end{infobox}

\newpage

% ============================================================
\section{Tab 5: Sectors}

Compara la performance de los distintos sectores del S\&P 500 usando ETFs sectoriales como proxies:

\begin{center}
\begin{tabular}{lll}
  \toprule
  \textbf{ETF} & \textbf{Sector} & \textbf{Color} \\
  \midrule
  XLK  & Tecnología               & Azul \\
  XLF  & Financiero               & Verde \\
  XLV  & Salud                    & Rojo \\
  XLE  & Energía                  & Naranja \\
  XLI  & Industrial               & Gris azulado \\
  XLY  & Consumo Discrecional     & Rosa \\
  XLP  & Consumo Básico           & Verde claro \\
  XLRE & Real Estate              & Marrón \\
  XLU  & Utilities                & Violeta \\
  XLB  & Materiales               & Verde oscuro \\
  XLC  & Comunicaciones           & Naranja oscuro \\
  \bottomrule
\end{tabular}
\end{center}

Todos los sectores se grafican en Base 100 para comparar su performance relativa en el período seleccionado. Incluye métricas básicas (Retorno, Volatilidad, Sharpe, Max DD) para cada sector.

\newpage

% ============================================================
\section{Datos y Fuentes}

\subsection{Yahoo Finance}

Todos los datos provienen de Yahoo Finance a través de requests HTTP directos:

\begin{itemize}
  \item \textbf{Precios históricos}: \api{query1.finance.yahoo.com/v8/finance/chart/\{TICKER\}}
  \item \textbf{Fundamentales}: \api{query1.finance.yahoo.com/v10/finance/quoteSummary/\{TICKER\}} con módulos \code{summaryDetail}, \code{defaultKeyStatistics}, \code{financialData}, \code{incomeStatementHistory}. Requiere autenticación por cookie (A3) + crumb.
  \item \textbf{Earnings}: mismo endpoint, módulo \code{earningsHistory}
  \item \textbf{Noticias}: \api{query1.finance.yahoo.com/v1/finance/search?q=\{TICKER\}\&newsCount=8}
\end{itemize}

\subsection{Autenticación Yahoo Finance (Crumb)}

Yahoo Finance requiere un par cookie+crumb para llamadas a quoteSummary:
\begin{enumerate}
  \item Visitar \code{https://fc.yahoo.com} → establece cookie \code{A3}
  \item Llamar \code{/v1/test/getcrumb} → devuelve el crumb como texto plano
  \item Incluir \code{?crumb=\{crumb\}} en cada request de quoteSummary
\end{enumerate}

El crumb se cachea globalmente (\code{\_YF\_CRUMB}) para reutilizarlo en todas las llamadas de la sesión.

\subsection{Cache y Performance}

\begin{center}
\begin{tabular}{ll}
  \toprule
  \textbf{Datos} & \textbf{Caché} \\
  \midrule
  Precios históricos  & Sin caché (request directo a Yahoo) \\
  Noticias            & 30 minutos en memoria (\code{\_NEWS\_CACHE}) \\
  Resumen cuantitativo & 5 minutos por ticker (\code{\_AI\_SUMMARY\_CACHE}) \\
  Fundamentales       & Sin caché (Yahoo puede fallar) \\
  Crumb Yahoo         & Toda la sesión (\code{\_YF\_CRUMB}) \\
  \bottomrule
\end{tabular}
\end{center}

\newpage

% ============================================================
\section{Guía de Uso Rápido}

\subsection{Comparison}
\begin{enumerate}
  \item Selecciona las empresas con los chips
  \item Define el rango de fechas
  \item Pulsa \textbf{Fetch Prices}
  \item Activa \textbf{30D Forecast} para ver proyecciones. Cambia entre GBM, GARCH y Merton
  \item Exporta a CSV con el botón \textbf{↓ CSV}
\end{enumerate}

\subsection{Technical}
\begin{enumerate}
  \item Selecciona el stock en el dropdown
  \item Define el rango de fechas y pulsa \textbf{Analyze}
  \item Activa/desactiva SMA20, SMA50, SMA200 y BB con los botones
  \item Hover sobre los triángulos \textbf{E} para ver datos de earnings
  \item Las señales automáticas aparecen bajo los gráficos
  \item Elige una estrategia y pulsa \textbf{$\triangleright$ Run Backtest}
\end{enumerate}

\subsection{Analytics}
\begin{enumerate}
  \item Selecciona las empresas, define fechas, pulsa \textbf{Run Analytics}
  \item Pulsa \textbf{Load Fundamentals} para activar las secciones de valuación
  \item \textbf{Precio Objetivo de Analistas}: revisa el consenso de Wall Street y el upside implícito
  \item \textbf{DCF Interactivo}: ajusta los sliders de crecimiento y WACC para ver el valor intrínseco estimado; compara con el precio actual
  \item El \textbf{Score de Inversión} combina Sharpe, Alpha, Calmar, Sortino, Drawdown y Beta en un puntaje 0--100
  \item Los \textbf{Insights automáticos} resumen los hallazgos clave del período
  \item El \textbf{Radar de Valuación} permite comparar todos los stocks en 6 dimensiones simultáneas
\end{enumerate}

\subsection{Portfolio}
\begin{enumerate}
  \item Selecciona las empresas y define los pesos (\%)
  \item Pulsa \textbf{Run Portfolio}
  \item Compara tu portafolio vs SPY y vs el Tangent Portfolio
\end{enumerate}

\subsection{Sectors}
\begin{enumerate}
  \item Define el rango de fechas y pulsa \textbf{Fetch Sectors}
  \item Compara la performance relativa de los 11 sectores del S\&P 500
\end{enumerate}

\newpage

% ============================================================
\section{Glosario}

\begin{description}[style=nextline, leftmargin=3cm, font=\bfseries\color{gold}]
  \item[Alpha ($\alpha$)] Retorno de una acción por encima de lo que predice el CAPM dado su nivel de beta. Mide la habilidad de un activo para generar valor independiente del mercado.
  \item[Base 100] Normalización que fija el precio inicial en 100 para comparar retornos entre activos de distinto precio absoluto.
  \item[Beta ($\beta$)] Sensibilidad del retorno de un activo a los movimientos del mercado. $\beta > 1$ amplifica; $\beta < 1$ amortigua.
  \item[Bollinger Bands] Bandas estadísticas a $\pm 2\sigma$ alrededor de una SMA. Contienen el 95\% de los precios bajo normalidad.
  \item[Buy \& Hold (B\&H)] Estrategia pasiva de comprar y mantener sin operar. Referencia de comparación para cualquier estrategia activa.
  \item[Calmar Ratio] Retorno anualizado dividido por el máximo drawdown absoluto. Mide cuánto retorno se genera por cada unidad de caída máxima.
  \item[Death Cross] SMA50 cruzando bajo SMA200. Señal técnica bajista de largo plazo.
  \item[Drawdown] Caída porcentual desde el último máximo. Mide pérdidas en el tiempo para un inversor que compró en el pico.
  \item[EPS] Earnings Per Share. Ganancia neta de la empresa dividida por el número de acciones en circulación.
  \item[EV/EBITDA] Enterprise Value dividido por EBITDA. Múltiplo de valoración que permite comparar empresas con distinto apalancamiento.
  \item[GARCH] Generalized Autoregressive Conditional Heteroskedasticity. Modelo que captura la variabilidad temporal de la volatilidad.
  \item[GBM] Geometric Brownian Motion. Proceso estocástico clásico para modelar precios de acciones con drift y difusión constantes.
  \item[Golden Cross] SMA50 cruzando sobre SMA200. Señal técnica alcista de largo plazo.
  \item[MACD] Moving Average Convergence Divergence. Indicador de momentum basado en la diferencia entre EMAs de 12 y 26 días.
  \item[Merton Jump-Diffusion] Extensión del GBM con saltos de Poisson. Captura eventos extremos (crashes, rallies).
  \item[PEG Ratio] P/E dividido por la tasa de crecimiento de ganancias. PEG $< 1$ puede indicar subvaloración relativa.
  \item[RSI] Relative Strength Index. Oscilador de momentum (0--100). $> 70$ = sobrecompra; $< 30$ = sobreventa.
  \item[Sharpe Ratio] Retorno en exceso sobre la tasa libre de riesgo por unidad de volatilidad total.
  \item[SMA] Simple Moving Average. Promedio aritmético de los últimos $n$ precios de cierre.
  \item[Sortino Ratio] Variante del Sharpe que divide por la volatilidad bajista únicamente. Más justo para activos con retornos asimétricos positivos.
  \item[Tangent Portfolio] Portafolio de máximo Sharpe Ratio en la frontera eficiente de Markowitz. Punto donde la Línea del Mercado de Capitales es tangente a la frontera.
  \item[Vol-Adjusted] Retorno base-100 dividido por la volatilidad anualizada. Penaliza activos que generaron el mismo retorno con mayor riesgo.
  \item[DCF] Discounted Cash Flow. Método de valoración que estima el valor intrínseco de una empresa descontando sus flujos de caja libres proyectados al presente usando el WACC.
  \item[WACC] Weighted Average Cost of Capital. Tasa de descuento que pondera el costo de la deuda y el costo del patrimonio según la estructura de capital. Aquí se aproxima como $r_f + \beta \times \text{ERP}$ con ERP = 5.5\%.
  \item[Free Cash Flow (FCF)] Flujo de caja libre. Caja generada por las operaciones del negocio después de capex. Es la base del modelo DCF.
  \item[Valor Terminal (TV)] Valor presente de todos los flujos de caja más allá del período explícito de proyección (año 10), estimado con el modelo de crecimiento perpetuo de Gordon.
  \item[ERP] Equity Risk Premium. Prima de riesgo del mercado accionario sobre la tasa libre de riesgo. Valor usado: 5.5\%.
  \item[Precio Objetivo de Analistas] Estimación del precio futuro de una acción publicada por analistas de sell-side. El consenso es la media de todas las estimaciones activas.
  \item[Radar de Valuación] Visualización tipo spider chart que compara múltiples activos en varias dimensiones simultáneas usando normalización min-max.
  \item[Min-Max Scaling] Normalización que transforma un valor al rango $[0, 100]$: $x' = (x - \min) / (\max - \min) \times 100$.
\end{description}

\vfill
\begin{center}
  \textcolor{gray}{\small MAG7 Tracker · Documentación Técnica v2.0 · Para actualizar: editar \texttt{docs/mag7\_doc.tex} y correr \texttt{docs/build.bat}}
\end{center}

\end{document}
